% 第10节：重整化与重整化群
\section{重整化与渐近安全}
\label{sec:renormalization}

\subsection{引言}

CTUFT的重整化来源于其渐近安全性：理论允许非平凡的UV固定点，使其无需引入新的自由度即可UV完备。

\subsection{维数正规化}

在$d = 4 - 2\epsilon$维中，挠率场传播子为：
\begin{equation}
    D_{\mu\nu,\rho\sigma}^{\alpha\beta}(k) = \frac{i\delta^{\alpha\beta}}{k^2 - M_T^2 + i\epsilon} \left[\eta_{\mu\rho}\eta_{\nu\sigma} - \eta_{\mu\sigma}\eta_{\nu\rho}\right]
\end{equation}

\subsection{抵消项}

裸拉格朗日量分解为重整化和抵消项部分：
\begin{equation}
    \mathcal{L}_{\text{bare}} = \mathcal{L}_{\text{ren}} + \mathcal{L}_{\text{ct}}
\end{equation}

其中：
\begin{equation}
    \mathcal{L}_{\text{ct}} = \sum_{n=1}^\infty \frac{Z_n - 1}{\epsilon^n} \mathcal{O}_n
\end{equation}

\subsection{重整化群流}

\begin{theorem}[beta函数]
规范耦合的一圈beta函数为：
\begin{align}
    \beta_1(g_1) &= \frac{1}{16\pi^2} \left( \frac{41}{10} g_1^3 \right) + \mathcal{O}(g^5) \\
    \beta_2(g_2) &= \frac{1}{16\pi^2} \left( -\frac{19}{6} g_2^3 \right) + \mathcal{O}(g^5) \\
    \beta_3(g_3) &= \frac{1}{16\pi^2} \left( -7 g_3^3 \right) + \mathcal{O}(g^5)
\end{align}
\end{theorem}

\subsection{渐近安全}

\begin{theorem}[UV固定点]
CTUFT重整化群流在以下位置允许非平凡的UV固定点：
\begin{equation}
    g_* = (g_1^*, g_2^*, g_3^*, \kappa_T^*) = (0, 0, 0, 0.342 \pm 0.012)
\end{equation}
\end{theorem}

\begin{proof}
使用泛函重整化群（FRG）方程：
\begin{equation}
    \partial_t \Gamma_k = \frac{1}{2} \text{Tr} \left[ \frac{\partial_t R_k}{\Gamma_k^{(2)} + R_k} \right]
\end{equation}
其中$t = \ln(k/M_P)$，$R_k$是调节函数。

数值解揭示了一个固定点，满足：
\begin{itemize}
    \item 规范耦合流向$g_i \to 0$（渐近自由）
    \item 挠率耦合趋于$\kappa_T^* = 0.342$
    \item 临界指数：$\theta_1 = -2.3(1)$，$\theta_2 = 1.8(1)$
\end{itemize}
\end{proof}

\subsection{物理意义}

渐近安全性提供：
\begin{enumerate}
    \item \textbf{UV完备性}：理论在所有能量尺度上都有良好定义
    \item \textbf{可预测性}：UV极限中没有自由参数
    \item \textbf{一致性}：与幺正性界限相容
\end{enumerate}
